\chapter{Functies van een reële veranderlijke}\label{ch:b2:realfun}

Voordat de analyse overgaat op functies van functies
(\cref{ch:b2:funcseq}) loont het het landschap in één veranderlijke
fijner te kennen dan bachelorjaar 1 vroeg: hoe discontinu een
monotone functie kan zijn, hoe regelmatig een convexe functie moet
zijn, en welke bijzondere eigenschappen afgeleiden genieten
(Darboux). Deze structurele resultaten zijn kort, scherp en geliefd
bij examinatoren.

\section{Monotone functies}

\begin{theorem}[Regulariteit van monotone functies]\label{thm:b2:realfun:monotone}
Zij $f \colon I \to \R$ stijgend op een interval.
\begin{enumerate}
  \item In elk inwendig punt $a$ bestaan de eenzijdige limieten:
        \[
        f(a^-) = \sup_{x < a} f(x) \;\leq\; f(a) \;\leq\;
        f(a^+) = \inf_{x > a} f(x) ;
        \]
        elke discontinuïteit is een \emph{sprong}.
  \item De verzameling van de discontinuïteiten van $f$ is
        hoogstens \omterm{def:b2:structures:countable}{aftelbaar}.\index{monotone functie!discontinuïteiten}
\end{enumerate}
\end{theorem}

\begin{proof}
(1) De verzameling $\{f(x) : x < a\}$ is niet leeg en van boven
door $f(a)$ begrensd: haar supremum $s$ voldoet aan $f(x) \to s$
als $x \to a^-$ (bij gegeven $\varepsilon$ is er een $f(x_0) > s -
\varepsilon$, en de monotonie sluit $f(x) \in \intoc{s -
\varepsilon}{s}$ in voor $x \in \intoo{x_0}{a}$). Symmetrisch aan
de rechterkant.

(2) Hecht aan elke discontinuïteit $a$ het niet-lege \omterm{def:b2:metric:topology}{open} interval
$J_a = \intoo{f(a^-)}{f(a^+)}$ (een echte sprong). Voor
discontinuïteiten $a < b$ zijn $J_a$ en $J_b$ disjunct: voor elke
$c$ ertussen geldt $f(a^+) \leq f(c) \leq f(b^-)$. Elke $J_a$ bevat
een rationaal getal, en verschillende discontinuïteiten krijgen
verschillende rationale getallen: dat is een injectie van de
verzameling discontinuïteiten in $\Q$, dat \omterm{def:b2:structures:countable}{aftelbaar} is
(\cref{prop:b2:structures:countablestable}).
\end{proof}

\begin{example}
De grens is scherp: leg een opsomming $(r_n)$ van $\Q \cap
\intoo{0}{1}$ vast en zet $f(x) = \sum_{n : r_n \leq x} 2^{-n}$
(een definitie via een \omterm{def:b2:series:summable}{sommeerbare familie},
\cref{def:b2:series:summable}). Dan is $f$ stijgend op
$\intcc{0}{1}$ en precies in elk rationaal getal van
$\intoo{0}{1}$ discontinu (met sprong $2^{-n}$ in $r_n$): een
monotone functie \emph{kan} dus op een dichte \omterm{def:b2:structures:countable}{aftelbare
verzameling} discontinu zijn.
\end{example}

\begin{example}[De sprongen kunnen de stijging niet overtreffen]\label{ex:b2:realfun:jumpbudget}
Voor stijgende $f$ op $\intcc{a}{b}$ hebben de sprongen een budget:
zijn $a < c_1 < \dots < c_m < b$ discontinuïteiten met sprongen
$s_i = f(c_i^+) - f(c_i^-) > 0$, kies dan tussenliggende punten $a
< c_1 < t_1 < c_2 < \dots$ en gebruik de monotonie op elk stuk:
\[
\sum_{i=1}^{m} s_i \;\leq\; f(b) - f(a) :
\]
de totale stijging begrenst het totale springen. Gevolg: voor elke
$k$ zijn er hoogstens $k\,\bigl(f(b) - f(a)\bigr)$
discontinuïteiten met sprong $\geq \frac1k$ --- een kwantitatieve
verfijning van \cref{thm:b2:realfun:monotone}~(2), want de
verzameling discontinuïteiten is de \omterm{def:b2:structures:countable}{aftelbare} vereniging over $k$
van deze eindige verzamelingen. Bij de functie met rationale
sprongen hierboven wordt het budget precies opgemaakt: de sprongen
$2^{-n}$ tellen op tot $1 = f(1^+) - f(0^-)$ in de voor de hand
liggende uitgebreide zin. Monotone functies mogen dicht springen,
maar alleen op een strikte toelage.
\end{example}

\section{Convexe functies}

\begin{lemma}[Hellingsongelijkheid]\label{lem:b2:realfun:slopes}
Zij $f$ convex op $I$ en zijn $x < y < z$ in $I$. Dan is
\[
\frac{f(y) - f(x)}{y - x}
\;\leq\; \frac{f(z) - f(x)}{z - x}
\;\leq\; \frac{f(z) - f(y)}{z - y} :
\]
de hellingen van de koorden stijgen in beide eindpunten.\index{convexe
functie!hellingsongelijkheid}
\end{lemma}

\begin{proof}
Schrijf $y = \frac{z - y}{z - x}\,x + \frac{y - x}{z - x}\,z$: een
convexe combinatie, want de twee coëfficiënten zijn positief en
tellen op tot $1$. De convexiteit geeft
\[
f(y) \;\leq\; \frac{z-y}{z-x}\,f(x) + \frac{y-x}{z-x}\,f(z).
\]
Trek voor de linkerongelijkheid $f(x)$ van beide leden af, met
$\frac{z-y}{z-x} - 1 = -\frac{y-x}{z-x}$:
\[
f(y) - f(x) \leq \frac{y - x}{z - x}\bigl(f(z) - f(x)\bigr),
\]
en deel door $y - x > 0$. Trek voor de rechterongelijkheid in
plaats daarvan af van $f(z)$:
\[
f(z) - f(y) \geq f(z) - \frac{z-y}{z-x}f(x) -
\frac{y-x}{z-x}f(z)
= \frac{z - y}{z - x}\bigl(f(z) - f(x)\bigr),
\]
en deel door $z - y > 0$. Beide stappen zijn dezelfde barycentrische
identiteit, gelezen tegen een ander eindpunt.
\end{proof}

\begin{theorem}[Regulariteit van convexe functies]\label{thm:b2:realfun:convexreg}
Zij $f$ convex op een interval $I$.
\begin{enumerate}
  \item In elk inwendig punt heeft $f$ eindige eenzijdige
        afgeleiden $f'_g \leq f'_d$; beide zijn stijgende functies
        van het punt; in het bijzonder is $f$ \omterm{def:b2:metric:continuity}{continu} op het
        inwendige van $I$ (maar mogelijk niet in de eindpunten).
  \item $f$ ligt boven elke \emph{steunlijn}: voor inwendige $a$ en
        elke $m \in \intcc{f'_g(a)}{f'_d(a)}$ geldt
        \[
        f(x) \geq f(a) + m(x - a) \qquad (x \in I).
        \]
  \item (Jensen, met gewichten) Voor $x_i \in I$ en gewichten
        $\lambda_i \geq 0$ met $\sum\lambda_i = 1$ geldt
        \[
        f\Bigl(\sum_i \lambda_i x_i\Bigr) \leq \sum_i \lambda_i
        f(x_i) .\index{ongelijkheid van Jensen}
        \]
\end{enumerate}
\end{theorem}

\begin{proof}
(1) Leg een inwendig punt $a$ vast. Volgens
\cref{lem:b2:realfun:slopes} is de helling $\tau(h) = \frac{f(a +
h) - f(a)}{h}$ een stijgende functie van $h$ (aan beide kanten, en
$\tau(h_-) \leq \tau(h_+)$ voor $h_- < 0 < h_+$). Bijgevolg heeft
$\tau$ een eindige limiet als $h \to 0^-$ (stijgend, van boven
begrensd door elke rechterhelling) --- dat is $f'_g(a)$ --- en als
$h \to 0^+$ ($f'_d(a)$), met $f'_g(a) \leq f'_d(a)$. Eindige
eenzijdige afgeleiden dwingen de \omterm{def:b2:metric:continuity}{continuïteit} in $a$ af. Monotonie
in het punt: voor inwendige $a < b$ is $f'_d(a) \leq \frac{f(b) -
f(a)}{b - a} \leq f'_g(b)$, opnieuw volgens de
hellingsongelijkheid.

(2) Voor $x > a$: $\frac{f(x) - f(a)}{x - a} \geq f'_d(a) \geq m$;
voor $x < a$: $\frac{f(a) - f(x)}{a - x} \leq f'_g(a) \leq m$.
Beide laten zich tot de bewering herschikken.

(3) Inductie naar het aantal punten, precies als in het volume van
bachelorjaar 1 (het geval van twee punten is de definitie) --- of
in één klap: pas (2) toe in $a = \sum\lambda_i x_i$ en middel de
steunlijnongelijkheden in de punten $x_i$ met de gewichten
$\lambda_i$: $\sum_i \lambda_i f(x_i) \geq f(a) +
m\sum_i\lambda_i(x_i - a) = f(a)$.
\end{proof}

\begin{omfigure}
\begin{tikzpicture}
\begin{axis}[omaxis, width=8.2cm, height=5.4cm,
    xmin=-2.4, xmax=2.5, ymin=-0.9, ymax=4.5,
    xtick={-1.5, 0.5, 2}, ytick=\empty]
  \addplot[omDef, very thick, domain=-2.1:2.15, samples=80]
    {x^2};
  \addplot[omThm, thick, domain=-1.5:2, samples=2]
    {0.5*x + 3};
  \addplot[omProp, thick, domain=-1.1:2.2, samples=2]
    {x - 0.25};
  \node[font=\small, omThm] at (axis cs:-1.05,3.15)
    {koorde};
  \node[font=\small, omProp] at (axis cs:1.75,0.95)
    {steunlijn};
  \node[font=\small, omDef] at (axis cs:-1.95,2.6) {$f$};
\end{axis}
\end{tikzpicture}

{\small Convexiteit in één plaatje: tussen $-1.5$ en $2$ blijft de
grafiek van $f(x) = x^2$ onder haar koorde (de definitie) en boven
de steunlijn in $x = 0.5$
(\cref{thm:b2:realfun:convexreg}~(2)) --- elke ongelijkheid uit de
weekendopgave van dit hoofdstuk is een herschikking van deze twee
posities.}
\end{omfigure}

\begin{example}[Discontinuïteit in een eindpunt]
Op $\intcc{0}{1}$ is de functie met $f(0) = 1$ en $f(x) = 0$ voor
$x > 0$ convex maar discontinu in het eindpunt $0$: uitspraak (1)
is dus scherp.
\end{example}

\begin{example}[Knikken en de waaier van steunlijnen]\label{ex:b2:realfun:corner}
Voor $f(x) = \abs x$ in $a = 0$ zijn de eenzijdige afgeleiden
$f'_g(0) = -1$ en $f'_d(0) = +1$, en
\cref{thm:b2:realfun:convexreg}~(2) deelt een steunlijn uit voor
\emph{elke} helling $m \in \intcc{-1}{1}$:
\[
\abs x \geq m\,x \qquad (x \in \R,\ -1 \leq m \leq 1),
\]
telkens met gelijkheid precies op een halfrechte of in $0$. Een
convexe functie is in $a$ differentieerbaar precies wanneer de
waaier tot één enkele lijn inklapt ($f'_g(a) = f'_d(a)$); knikken
dragen een heel interval raaklijnen. Deze waaier is de
eindigdimensionale kiem van het \emph{subdifferentiaal} uit de
convexe optimalisatie --- en de reden dat convexe functies zo
robuust zijn: zelfs waar de afgeleide faalt, overleeft de
steunende meetkunde, en meer had het bewijs van Jensen niet
nodig.
\end{example}

\begin{example}[Ongelijkheid van de machtgemiddelden]\label{ex:b2:realfun:powermeans}
Voor $0 < p < q$ en positieve $x_i$ met gewichten $\lambda_i$ die
tot $1$ optellen geeft Jensen, toegepast op de convexe $t \mapsto
t^{q/p}$ in de punten $x_i^p$:
\[
\Bigl(\sum \lambda_i x_i^{p}\Bigr)^{1/p}
\leq \Bigl(\sum \lambda_i x_i^{q}\Bigr)^{1/q} :
\]
de machtgemiddelden stijgen met de exponent --- waaronder de
ongelijkheid tussen rekenkundig en kwadratisch gemiddelde, en, in
de limiet $p \to 0$ (\cref{exo:b2:realfun:6}), opnieuw de
ongelijkheid tussen rekenkundig en meetkundig gemiddelde.
\end{example}

\begin{omfigure}
\begin{tikzpicture}
\begin{axis}[omaxis, width=8.6cm, height=5.6cm,
    xmin=-8.5, xmax=8.5, ymin=0.5, ymax=4.5,
    xtick={-4,-1,1,2,4}, ytick={1,2,4},
    xlabel={$p$}, ylabel={$M_p$}]
  \addplot[omDef, very thick, domain=-8:-0.4, samples=90]
    {((1/3)*(1^x + 2^x + 4^x))^(1/x)};
  \addplot[omDef, very thick, domain=0.4:8, samples=90]
    {((1/3)*(1^x + 2^x + 4^x))^(1/x)};
  \addplot[omThm, dashed, domain=-8.5:8.5, samples=2] {4};
  \addplot[omThm, dashed, domain=-8.5:8.5, samples=2] {1};
  \node[font=\small, omDef] at (axis cs:-5.4,1.55)
    {$M_p(1,2,4)$};
  \node[font=\small, omThm] at (axis cs:-6.4,3.75)
    {$\max = 4$};
  \node[font=\small, omThm] at (axis cs:6.3,1.25)
    {$\min = 1$};
  \node[font=\small, omProp] at (axis cs:0.2,1.83) {$2$};
\end{axis}
\end{tikzpicture}

{\small Het machtgemiddelde $M_p$ van de waarden $1, 2, 4$ (gelijke
gewichten), als functie van de exponent $p$: stijgend van $\min =
1$ (als $p \to -\infty$) tot $\max = 4$ (als $p \to +\infty$), via
het harmonische ($p = -1$), het meetkundige (het gat bij $p = 0$,
waarde $2$), het rekenkundige ($p = 1$) en het kwadratische ($p =
2$) gemiddelde. De hele keten van klassieke ongelijkheden tussen
gemiddelden is één stijgende kromme --- bewezen in Deel III van de
weekendopgave van dit hoofdstuk.}
\end{omfigure}

\begin{example}[Maximale entropie]\label{ex:b2:realfun:entropy}
Voor een kansvector $(p_1, \dots, p_n)$ (positief, met som $1$)
voldoet de entropie $H(p) = -\sum_i p_i\ln p_i$ aan
\[
H(p) \leq \ln n ,
\qquad\text{met gelijkheid dan en slechts dan als } p_i =
\frac1n \text{ voor alle } i .
\]
Bewijs met Jensen (\cref{thm:b2:realfun:convexreg}~(3)) toegepast
op de \emph{concave} $\ln$ met gewichten $p_i$ in de punten
$\frac{1}{p_i}$:
\[
H(p) = \sum_i p_i\ln\frac{1}{p_i}
\leq \ln\Bigl(\sum_i p_i\,\frac1{p_i}\Bigr) = \ln n ,
\]
waarbij gelijkheid alle punten $\frac1{p_i}$ gelijk afdwingt
(strikte concaviteit), dat wil zeggen: $p$ uniform. Gelijkwaardig
is dit \cref{exo:b2:realfun:7} met uniforme $q$. Onzekerheid wordt
gemaximaliseerd door onwetendheid die uniform is uitgesmeerd ---
het variationele principe achter codering, statistische mechanica
en de entropie zoals zij in \cref{ch:b2:randomvar} opduikt.
\end{example}

\begin{method}[De convexe functie achter een ongelijkheid vinden]\label{met:b2:realfun:spotting}
De meeste klassieke ongelijkheden zijn Jensen in vermomming; om er
een te ontkleden: (1) normaliseer zo dat er een \emph{gewogen
gemiddelde} verschijnt (positieve gewichten met som $1$ --- deel
zo nodig door een totale massa); (2) kijk welke functie binnen en
welke buiten het gemiddelde wordt toegepast: de bewering
``$f(\text{gemiddelde}) \leq$ gemiddelde van $f$'' noemt de convexe
$f$; (3) certificeer de convexiteit met de tweede afgeleide en
behandel de gelijkheid met de striktheid; (4) is er geen gemiddelde
zichtbaar, neem dan eerst logaritmen --- producten en machten
worden gemiddelden, en de concaviteit van $\ln$ draagt de
ongelijkheid tussen rekenkundig en meetkundig gemiddelde, die van
Young en hun verwanten (de weekendopgave van dit hoofdstuk doorloopt
de stappen 1--4 op elk ervan). Onthullen zelfs logaritmen geen
gemiddelde, lees de ongelijkheid dan als monotonie van hellingen
(\cref{lem:b2:realfun:slopes}) --- uitspraken over
superadditiviteit zoals \cref{exo:b2:realfun:9} horen daar
thuis.
\end{method}

\begin{remark}[Klassieke valkuilen]\label{rem:b2:realfun:pitfalls}
(i) Convexiteit blijft niet onder producten behouden: $x$ en $(x -
1)^2$ zijn convex op $\intcc{0}{2}$, maar hun product $x(x-1)^2$
heeft tweede afgeleide $6x - 4$, negatief op $\intco{0}{\frac23}$
--- dus niet convex; en zonder monotonie blijft convexiteit
evenmin onder samenstelling behouden (\cref{exo:b2:realfun:10}).
(ii) Jensen keert om voor concave functies: de helft van de
klassieke ongelijkheden is de concave $\ln$-versie; de convexe vorm
op $\ln$ toepassen is de snelste manier om de ongelijkheid tussen
rekenkundig en meetkundig gemiddelde \emph{achterstevoren} te
bewijzen. (iii) Convexiteit in het midden alleen impliceert geen
convexiteit --- daarvoor is \omterm{def:b2:metric:continuity}{continuïteit} (of alleen al
begrensdheid) nodig (\cref{exo:b2:realfun:8}); de pathologische
tegenvoorbeelden liggen buiten de axioma's van dit boek. (iv) Een
convexe functie op een \emph{\omterm{def:b2:metric:topology}{open}} interval is \omterm{def:b2:metric:continuity}{continu}, zelfs
lokaal \omterm{def:b2:metric:continuity}{Lipschitz} (\cref{exo:b2:realfun:12}); in de eindpunten is
niets gratis. (v) Afgeleiden gehoorzamen aan Darboux maar hoeven
niet \omterm{def:b2:metric:continuity}{continu} te zijn (\cref{ex:b2:realfun:oscillation}): ``$f'$
heeft geen sprongen'' betekent nooit ``$f'$ is \omterm{def:b2:metric:continuity}{continu}''.
\end{remark}

\section{De eigenschap van Darboux}

\begin{theorem}[Darboux]\label{thm:b2:realfun:darboux}
Zij $f$ differentieerbaar op een interval $I$. Dan neemt $f'$ elke
waarde tussen twee van haar waarden aan --- ook al hoeft $f'$ niet
\omterm{def:b2:metric:continuity}{continu} te zijn.\index{stelling van Darboux}
\end{theorem}

\begin{proof}
Zijn $a < b$ in $I$ en ligt $v$ strikt tussen $f'(a)$ en $f'(b)$,
zeg $f'(a) < v < f'(b)$. De functie $g(x) = f(x) - vx$ is
differentieerbaar met $g'(a) < 0 < g'(b)$: haar minimum op
$\intcc{a}{b}$ (aangenomen: \omterm{def:b2:metric:continuity}{continuïteit} op een \omterm{def:b2:metric:compact}{compacte}
verzameling) ligt niet in $a$ (net na $a$ daalt $g$ onder $g(a)$)
en evenmin in $b$ (net vóór $b$ ligt $g$ onder $g(b)$): het is dus
inwendig, en daar is $g'(c) = 0$, dat wil zeggen $f'(c) = v$. (Dit
was een oefening met sterren in bachelorjaar 1; haar plaats in de
theorie is hier.)
\end{proof}

\begin{example}[Welke functies zijn afgeleiden?]\label{ex:b2:realfun:whichderivatives}
De stelling van Darboux is een niet-bestaansmachine. De
afrondingsfunctie $\lfloor x\rfloor$ is op $\R$ de afgeleide van
geen enkele functie: zij neemt de waarden $0$ en $1$ aan maar slaat
$\frac12$ op $\intcc{0}{1}$ over, en dat verbiedt
\cref{thm:b2:realfun:darboux} voor afgeleiden. Hetzelfde oordeel
treft elke functie met een sprong --- het teken, de
heavisidefunctie, alle trapfuncties --- hoe onschuldig zij er ook
uitzien; hun ``primitieven'' ($\abs x$ voor het teken, enzovoort)
bestaan alleen buiten de sprong en knopen daar met een knik aaneen.
Daartegenover \emph{is} de wild discontinue $f'$ uit
\cref{ex:b2:realfun:oscillation} wél een afgeleide --- haar
discontinuïteit is een oscillatie, en die verdraagt Darboux. De
grens tussen beide gedragingen is precies het gevolg over
sprongen hieronder.
\end{example}

\begin{corollary}\label{cor:b2:realfun:nojumps}
Een afgeleide heeft geen sprongdiscontinuïteiten: bestaan $f'(a^-)$
en $f'(a^+)$, dan zijn zij gelijk aan $f'(a)$. De
discontinuïteiten van een afgeleide zijn altijd van het
oscillatietype (de afgeleide van $x^2\sin\frac1x$ in $0$, volume
van bachelorjaar 1).
\end{corollary}

\begin{proof}
Bestaat $f'(a^+) = \lim_{x\to a^+} f'(x)$ en verschilt zij van
$f'(a)$, dan zouden waarden er strikt tussenin door $f'$ op een
rechteromgeving worden overgeslagen --- in strijd met Darboux op de
intervallen $\intcc{a}{a + h}$. (Anders gezegd: de
middelwaardestelling dwingt $f'(a) = \lim_{h\to0^+}
\frac{f(a+h)-f(a)}{h} = f'(a^+)$ af, want het differentiequotiënt
is een waarde van $f'$ in een tussenliggend punt.) Aan de
linkerkant hetzelfde.
\end{proof}

\begin{example}[De canonieke oscillerende afgeleide]\label{ex:b2:realfun:oscillation}
Zij $f(x) = x^2\sin\frac1x$ voor $x \neq 0$ en $f(0) = 0$. In $0$:
$\bigl|\frac{f(h) - f(0)}{h}\bigr| = \abs{h\sin\frac1h} \leq
\abs h \to 0$, dus bestaat $f'(0) = 0$. Buiten $0$ is
\[
f'(x) = 2x\sin\frac1x - \cos\frac1x ,
\]
waarvan de eerste term naar $0$ gaat terwijl $\cos\frac1x$ op elk
interval $\intoo{0}{\delta}$ het hele $\intcc{-1}{1}$ doorloopt: de
limiet $f'(0^+)$ bestaat dus niet. Bijgevolg is $f'$ overal
gedefinieerd maar discontinu in $0$ --- en, precies zoals
\cref{cor:b2:realfun:nojumps} voorspelt, is de discontinuïteit een
oscillatie en geen sprong: op elke $\intoo{0}{\delta}$ veegt $f'$
nog steeds een heel interval rond $0$ af. Afgeleiden mogen wild
zijn, maar alleen op de manier die met Darboux verenigbaar is.
\end{example}

\begin{remark}[Waar dit hoofdstuk wordt gebruikt]
De convexiteit is de motor van de ongelijkhedenindustrie: de
weekendopgave van dit hoofdstuk maakt er Young, Hölder, Minkowski
en de keten van machtgemiddelden mee, die de normentheorie van
\cref{ch:b2:nvs} en de integraalschattingen van
\cref{ch:b2:integration} weer verbruiken; Jensen duikt in de
kansrekening opnieuw op als de momentongelijkheden van
\cref{ch:b2:randomvar}. De regulariteit van monotone functies keert
terug in \cref{ch:b2:integration} (monotone functies zijn
integreerbaar) en, in het volume van bachelorjaar 3, als de
differentieerbaarheid bijna overal van monotone functies --- waar
``\omterm{def:b2:structures:countable}{aftelbaar} veel sprongen'' de eerste stap van de theorie van
Lebesgue wordt.
\end{remark}

\section{Oefeningen}

\begin{exercise}[$\star$]\label{exo:b2:realfun:1}
Bepaal de verzamelingen van discontinuïteiten en de sprongen van
$\lfloor x \rfloor$, $\;x - \lfloor x\rfloor$, $\;\lfloor x \rfloor
+ \sqrt{x - \lfloor x\rfloor}$, en van de functie uit het voorbeeld
na \cref{thm:b2:realfun:monotone}, beperkt tot de dyadische
rationale getallen $r_n$.
\end{exercise}

\begin{exercise}[$\star$]\label{exo:b2:realfun:2}
Bewijs dat een stijgende functie $f \colon I \to \R$ met de
tussenwaarde-eigenschap (het beeld van elk deelinterval is een
interval) \omterm{def:b2:metric:continuity}{continu} is.
\end{exercise}

\begin{exercise}[$\star$]\label{exo:b2:realfun:3}
Welke van de volgende functies zijn convex op hun domein? $x
\mapsto x\ln x$ ($x > 0$); $\;x \mapsto \ln(1 + \eu^x)$; $\;x
\mapsto \sqrt{1 + x^2}$; $\;x \mapsto x^3$.
\end{exercise}

\begin{exercise}[$\star\star$]\label{exo:b2:realfun:4}
Zij $f$ convex op $\R$ en van boven begrensd. Bewijs dat $f$
constant is. \emph{(Is $f(a) \neq f(b)$, dan plant de
hellingsongelijkheid de koordehelling ongelijk aan nul voort:
voorbij het punt met de grootste waarde groeit $f$ minstens lineair
--- in strijd met de begrensdheid. Behandel beide tekens van de
helling.)} Leid af dat een convexe functie op $\R$ met aan beide
zijden een asymptoot affien is.
\end{exercise}

\begin{exercise}[$\star\star$]\label{exo:b2:realfun:5}
Zij $f$ differentieerbaar op $I$ met $f'$ monotoon. Bewijs dat $f'$
\omterm{def:b2:metric:continuity}{continu} is \emph{(combineer \cref{thm:b2:realfun:monotone} en
\cref{cor:b2:realfun:nojumps})}.
\end{exercise}

\begin{exercise}[$\star\star$]\label{exo:b2:realfun:6}
(Het meetkundige gemiddelde als limiet) Bewijs voor positieve $x_i$
en gewichten $\lambda_i$ met som $1$ dat
\[
\lim_{p \to 0^+} \Bigl(\sum_i \lambda_i x_i^p\Bigr)^{1/p}
= \prod_i x_i^{\lambda_i} ,
\]
via $x_i^p = \eu^{p\ln x_i} = 1 + p\ln x_i + O(p^2)$, en leid
daaruit met \cref{ex:b2:realfun:powermeans} de gewogen
ongelijkheid tussen rekenkundig en meetkundig gemiddelde af.
\end{exercise}

\begin{exercise}[$\star\star$]\label{exo:b2:realfun:7}
(Entropieongelijkheid) Bewijs met de strikte convexiteit van $t
\mapsto t\ln t$ dat voor positieve $p_i, q_i$ met $\sum p_i = \sum
q_i = 1$ geldt
\[
\sum_i p_i \ln\frac{p_i}{q_i} \geq 0 ,
\]
met gelijkheid dan en slechts dan als $p = q$. \emph{(Schrijf het
linkerlid als $\sum q_i\, \varphi\bigl(\frac{p_i}{q_i}\bigr)$ met
$\varphi(t) = t\ln t$ en pas Jensen toe met gewichten $q_i$.)}
\end{exercise}

\begin{exercise}[$\star\star\star$]\label{exo:b2:realfun:8}
(Convexiteit in het midden) $f \colon I \to \R$ heet \emph{convex
in het midden} wanneer altijd $f\bigl(\frac{x+y}{2}\bigr) \leq
\frac{f(x) + f(y)}{2}$. Bewijs dat een \emph{\omterm{def:b2:metric:continuity}{continue}} functie die
convex in het midden is, convex is. \emph{(Stel de
convexiteitsongelijkheid voor dyadische gewichten $\frac{k}{2^m}$
vast met inductie naar $m$, en ga daarna met de dichtheid en de
\omterm{def:b2:metric:continuity}{continuïteit} naar de limiet.)}
\end{exercise}

\begin{exercise}[$\star\star\star$]\label{exo:b2:realfun:9}
Zij $f$ convex op $\intco{0}{+\infty}$ met $f(0) \leq 0$. Bewijs
dat $x \mapsto \frac{f(x)}{x}$ stijgt op $\intoo{0}{+\infty}$, en
leid af dat voor convexe $f$ met $f(0) = 0$ geldt: $f(x + y) \geq
f(x) + f(y)$ voor $x, y \geq 0$ (superadditiviteit).
\end{exercise}

\begin{exercise}[$\star$]\label{exo:b2:realfun:10}
Zij $f$ convex op $I$ en $g$ convex en \emph{stijgend} op een
interval dat $f(I)$ bevat. Bewijs dat $g \circ f$ convex is, en
toon met een tegenvoorbeeld aan dat de monotonie van $g$ niet mag
vervallen.
\end{exercise}

\begin{exercise}[$\star\star$]\label{exo:b2:realfun:11}
(Hermite--Hadamard) Zij $f$ convex en \omterm{def:b2:metric:continuity}{continu} op $\intcc{a}{b}$.
Bewijs dat
\[
f\Bigl(\frac{a+b}{2}\Bigr) \;\leq\; \frac{1}{b -
a}\int_a^b f(t)\,\dd t \;\leq\; \frac{f(a) + f(b)}{2} .
\]
\emph{(Links: integreer een steunlijn in het midden. Rechts:
begrens $f$ door de koorde.)}
\end{exercise}

\begin{exercise}[$\star\star\star$]\label{exo:b2:realfun:12}
Bewijs dat een convexe functie op een \emph{\omterm{def:b2:metric:topology}{open}} interval $I$
lokaal \omterm{def:b2:metric:continuity}{Lipschitz} is: voor elk segment $\intcc{a}{b} \subseteq I$ en
elke marge $\delta > 0$ met $\intcc{a - \delta}{b + \delta}
\subseteq I$ is de beperking van $f$ tot $\intcc{a}{b}$ \omterm{def:b2:metric:continuity}{Lipschitz},
met constante $\max\Bigl(\bigl|\frac{f(a) - f(a -
\delta)}{\delta}\bigr|, \bigl|\frac{f(b + \delta) -
f(b)}{\delta}\bigr|\Bigr)$ \emph{(sluit elke koordehelling met de
hellingsongelijkheid tussen deze twee in)}.
\end{exercise}

\section{Probleem: De gereedschapskist van de convexiteit}

Eén definitie --- de koorde boven de grafiek --- brengt de hele
gereedschapskist van de klassieke ongelijkheden voort. Deze
weekendopgave bouwt haar in logische volgorde op:
convexiteitscriteria en strikte Jensen, dan Young, Hölder en
Minkowski (de geboorteakten van de $p$-normen), de volledige keten
van machtgemiddelden van het minimum tot het maximum, en twee
kroonopbrengsten --- de ongelijkheid van Carleman, en Hölder
gelezen als een dualiteit. Alles wordt bewezen; niets wordt
ingevoerd.

\begin{problem}[{Weekendopgave --- Young, Hölder, Minkowski en de
keten van machtgemiddelden}]
\label{pb:b2:realfun:1}
Overal zijn $p, q > 1$ \emph{toegevoegde exponenten}: $\frac1p +
\frac1q = 1$; vectoren zijn $a = (a_1, \dots, a_n) \in \R^n$; en
gewichten $\lambda_i > 0$ voldoen aan $\sum_i\lambda_i = 1$.

\medskip
\noindent\textbf{Deel I --- Criteria en strikte Jensen.}
\begin{enumerate}
  \item Zij $f$ differentieerbaar op een interval $I$. Bewijs dat
        $f$ convex is dan en slechts dan als $f'$ stijgt \emph{(de
        ene richting door in de hellingsongelijkheid
        \cref{lem:b2:realfun:slopes} naar de limiet te gaan, de
        andere met de middelwaardestelling)}. Leid het
        $C^2$-criterium $f'' \geq 0$ af.
  \item Neem aan dat $f'' > 0$ op $I$. Bewijs dat $f$
        \emph{strikt} convex is (strikte ongelijkheid voor $x \neq
        y$ en $\lambda \in \intoo01$), en dat een strikt convexe
        functie aan de ongelijkheid van Jensen
        (\cref{thm:b2:realfun:convexreg}~(3)) \emph{alleen} met
        gelijkheid voldoet wanneer alle $x_i$ samenvallen.
  \item Certificeer de grondstoffen van de gereedschapskist:
        $-\ln$ is strikt convex op $\intoo{0}{+\infty}$; $t \mapsto
        t^r$ is daar strikt convex voor $r > 1$ en strikt concaaf
        voor $0 < r < 1$; en $\exp$ is strikt convex op $\R$.
  \item (Ongelijkheid van Young) Bewijs voor $a, b \geq 0$ dat
        \[
        ab \;\leq\; \frac{a^p}{p} + \frac{b^q}{q},
        \]
        met gelijkheid dan en slechts dan als $a^p = b^q$
        \emph{(pas de concaviteit van $\ln$ toe op de twee punten
        $a^p, b^q$ met gewichten $\frac1p, \frac1q$)}.
  \item Leid uit de concaviteit van $\ln$ in één regel opnieuw de
        gewogen ongelijkheid tussen rekenkundig en meetkundig
        gemiddelde af:
        \[
        \prod_i x_i^{\lambda_i} \leq \sum_i\lambda_ix_i \qquad
        (x_i > 0),
        \]
        met het geval van gelijkheid; vergelijk dit met de weg via
        de limiet uit \cref{exo:b2:realfun:6}.
\end{enumerate}

\medskip
\noindent\textbf{Deel II --- Hölder en Minkowski.} Schrijf
$\norm{a}_p = \bigl(\sum_i \abs{a_i}^p\bigr)^{1/p}$ en
$\norm{a}_\infty = \max_i\abs{a_i}$.
\begin{enumerate}[resume]
  \item (Hölder) Bewijs dat
        \[
        \sum_{i=1}^{n}\abs{a_ib_i} \;\leq\;
        \norm a_p\,\norm b_q ,
        \]
        met gelijkheid dan en slechts dan als de vectoren
        $(\abs{a_i}^p)$ en $(\abs{b_i}^q)$ evenredig zijn
        \emph{(normaliseer $\norm a_p = \norm b_q = 1$ en pas
        Young term voor term toe)}.
  \item Benoem de bijzondere gevallen: $p = q = 2$
        (Cauchy--Schwarz), en het randpaar $(p, q) = (1, \infty)$:
        formuleer en bewijs $\sum\abs{a_ib_i} \leq \norm
        a_1\norm b_\infty$.
  \item (Minkowski) Bewijs voor $p \geq 1$ dat
        \[
        \norm{a + b}_p \leq \norm a_p + \norm b_p
        \]
        \emph{(schrijf $\abs{a_i + b_i}^p \leq \abs{a_i +
        b_i}^{p-1}(\abs{a_i} + \abs{b_i})$ en pas Hölder op elk
        product toe)}. Besluit: $\norm\cdot_p$ is voor elke $p \in
        \intco{1}{+\infty}$ een \omterm{def:b2:nvs:norm}{norm} op $\R^n$, wat het beeld van
        \cref{ch:b2:nvs} completeert.
  \item Integraalversies: formuleer en bewijs Hölder en Minkowski
        voor $\norm f_p = \bigl(\int_a^b\abs f^p\bigr)^{1/p}$ met
        $f, g$ \omterm{def:b2:metric:continuity}{continu} op $\intcc{a}{b}$ \emph{(dezelfde
        bewijzen, met de strikte positiviteit van de integraal
        voor de bespreking van de gelijkheid)}.
  \item Bewijs de monotonie $\norm a_q \leq \norm a_p$ voor $1
        \leq p \leq q$, de limiet $\norm a_p \to \norm a_\infty$
        als $p \to \infty$, en de omgekeerde vergelijking met de
        scherpe constante:
        \[
        \norm a_p \leq n^{\frac1p - \frac1q}\,\norm a_q
        \]
        \emph{(Hölder tegen de constante vector)}. Benoem de
        vectoren die elke gelijkheid realiseren.
  \item (Interpolatie) Bewijs voor $1 \leq p < r < q$ en $\theta
        \in \intoo01$ met $\frac1r = \frac\theta p +
        \frac{1-\theta}q$ dat
        \[
        \norm a_r \leq \norm a_p^{\theta}\,
        \norm a_q^{1-\theta}
        \]
        \emph{(pas Hölder met de exponenten $\frac{p}{\theta r}$
        en $\frac{q}{(1-\theta)r}$ toe op
        $\abs{a_i}^{\theta r}\abs{a_i}^{(1-\theta)r}$)}.
\end{enumerate}

\medskip
\noindent\textbf{Deel III --- De keten van machtgemiddelden,
\omterm{def:b2:metric:complete}{volledig}.} Zet voor $p \neq 0$ de waarde $M_p =
\bigl(\sum_i\lambda_i x_i^p\bigr)^{1/p}$ ($x_i > 0$), en $M_0 =
\prod_i x_i^{\lambda_i}$.
\begin{enumerate}[resume]
  \item Bewijs dat $p \mapsto M_p$ op heel $\R^*$ stijgt: behandel
        $p < q < 0$ met de identiteit $M_{-p}(x) =
        M_p(1/x)^{-1}$, en overbrug $0$ door aan te tonen dat $M_p
        \leq M_0 \leq M_q$ voor $p < 0 < q$ \emph{(pas de
        concaviteit van $\ln$ toe op $x_i^q$, en de omgekeerde
        ongelijkheid voor negatieve exponenten)}.
  \item Bewijs de limieten $M_p \to \max_i x_i$ als $p \to
        +\infty$ en $M_p \to \min_i x_i$ als $p \to -\infty$.
  \item Schrijf de keten $\min \leq \mathrm{HM} \leq \mathrm{GM}
        \leq \mathrm{AM} \leq \mathrm{QM} \leq \max$ uit voor
        gelijke gewichten, en bewijs het klassieke gevolg: voor
        positieve $a_1, \dots, a_n$ geldt
        \[
        \Bigl(\sum_i a_i\Bigr)\Bigl(\sum_i\frac1{a_i}\Bigr)
        \geq n^2 .
        \]
  \item Verbind gemiddelden met \omterm{def:b2:nvs:norm}{normen}: voor gelijke gewichten
        $\lambda_i = \frac1n$ is $M_p(x) = n^{-1/p}\norm x_p$.
        Verzoen in één zin over de factor $n^{-1/p}$ de twee
        monotonieën --- gemiddelden \emph{stijgen} met $p$ terwijl
        \omterm{def:b2:nvs:norm}{normen} \emph{dalen} (vraag 10).
  \item Bepaal de gevallen van gelijkheid langs de hele keten van
        vraag 14 (positieve gewichten): gelijkheid ergens dwingt
        af dat alle $x_i$ gelijk zijn --- de strikte convexiteit
        betaalt zich uit.
\end{enumerate}

\medskip
\noindent\textbf{Deel IV --- Opbrengsten.}
\begin{enumerate}[resume]
  \item (Young met een knop) Bewijs voor $a, b \geq 0$ en
        $\varepsilon > 0$ dat
        \[
        ab \leq \varepsilon\,\frac{a^p}{p} +
        \varepsilon^{-q/p}\,\frac{b^q}{q},
        \]
        en het werkpaardgeval $ab \leq \varepsilon a^2 +
        \frac{b^2}{4\varepsilon}$: de absorptietruc die in de hele
        analyse wordt gebruikt.
  \item (Op weg naar Carleman) Zij $c_k = \frac{(k+1)^k}{k^{k-1}}$.
        Bewijs de telescopische identiteit $\prod_{k=1}^{n}c_k =
        (n+1)^n$, en leid met de ongelijkheid tussen rekenkundig
        en meetkundig gemiddelde toegepast op de getallen $c_ka_k$
        af dat
        \[
        (a_1a_2\cdots a_n)^{1/n} \leq
        \frac{1}{n(n+1)}\sum_{k=1}^{n} c_k a_k
        \qquad (a_k > 0).
        \]
  \item (Ongelijkheid van Carleman) Sommeer over $n$, verwissel de
        sommatievolgorde (positieve \omterm{def:b2:series:summable}{sommeerbare families},
        \cref{thm:b2:series:fubini}), en gebruik $\sum_{n \geq
        k}\frac{1}{n(n+1)} = \frac1k$ en $c_k/k = \bigl(1 +
        \frac1k\bigr)^k < \eu$ om te besluiten: voor elke
        convergente $\sum a_k$ met positieve termen geldt
        \[
        \sum_{n=1}^{\infty}(a_1a_2\cdots a_n)^{1/n}
        \;\leq\; \eu\sum_{k=1}^{\infty}a_k .
        \]
  \item Bewijs voor $f$ \omterm{def:b2:metric:continuity}{continu} en positief op $\intcc{0}{1}$ dat
        \[
        \Bigl(\int_0^1 f\Bigr)\Bigl(\int_0^1\frac1f\Bigr)
        \geq 1,
        \]
        met gelijkheid dan en slechts dan als $f$ constant is
        \emph{(Cauchy--Schwarz op $\sqrt f\cdot\frac1{\sqrt
        f}$)}.
  \item (Meetkunde van de bollen) Toon met het geval van
        gelijkheid in Minkowski aan dat voor $1 < p < \infty$ de
        eenheidssfeer van $\norm\cdot_p$ geen enkel segment bevat
        (de \omterm{def:b2:nvs:norm}{norm} is strikt convex in de zin van
        \cref{pb:b2:nvs:1}), terwijl zij dat voor $p = 1$ en $p =
        \infty$ wel doet: geef de vlakke stukken.
\end{enumerate}

\medskip
\noindent\textbf{Deel V --- Dualiteit en synthese.}
\begin{enumerate}[resume]
  \item (Hölder als dualiteit) Bewijs dat voor elke $a \in \R^n$
        \[
        \norm a_p = \max_{\norm b_q \leq 1}\ \sum_i a_ib_i ,
        \]
        en geef een maximaliserende $b$ expliciet. (De $p$-norm is
        de \omterm{def:b2:linalg:dual}{duale} van de $q$-norm --- de eindigdimensionale kiem
        van de dualiteit van de $L^p$-ruimten.)
  \item (Momenten) Zij $X$ een stochastische variabele die eindig
        veel positieve waarden $x_i$ aanneemt met kansen
        $\lambda_i$. Herformuleer vraag 12 als: $r \mapsto
        \E[X^r]^{1/r}$ stijgt --- de momentongelijkheid (van
        Lyapunov), opnieuw te gebruiken in
        \cref{ch:b2:randomvar}.
  \item Los op met benoemd gereedschap, in telkens twee regels:
        (i) voor positieve $a, b, c$ geldt $a^3 + b^3 + c^3 \geq
        \frac{(a+b+c)^3}{9}$; (ii) voor positieve $x_1, \dots,
        x_n$ geldt $\bigl(\sum_i\sqrt{x_i}\bigr)^2 \leq n\sum_i
        x_i$.
  \item (Synthese) Teken de stamboom in vijf zinnen: van de
        definitie met de koorde naar het hellingslemma; van de
        hellingen via de steunlijnen naar Jensen; van de
        concaviteit van $\ln$ via Young en Hölder naar Minkowski
        en de $p$-normen; van Jensen via de keten van
        machtgemiddelden naar de momenten; van de ongelijkheid
        tussen rekenkundig en meetkundig gemiddelde naar Carleman.
        Noem de toppen (Hölder en Minkowski; Carleman) en zeg
        waar de gereedschapskist naartoe gaat: de
        $L^p$-ruimten van het volume van bachelorjaar 3, waarvan
        de axioma's precies de vragen 6 en 8 zijn.
\end{enumerate}
\end{problem}
